\section*{Bab \ref{ch:b3:orogenic-wedges} --- Sabuk Sesar Dorong dan Baji Orogen}

\begin{solution}{exo:b3:orogenic-wedges:1}
$\sin\ang{30} = 0.5$: $K = 2\times0.5/0.5 = 2$ dan $\sigma'_1/\sigma'_3 = 1 +
K = 3$.
\end{solution}

\begin{solution}{exo:b3:orogenic-wedges:2}
$\phi = \arctan0.6 = \ang{31.0}$, $\sin\phi = 0.514$, $K = 2.12$. Pada
dasar horizontal, $\alpha = 0.3/3.12 = \qty{0.0962}{rad} = \ang{5.5}$. Pada
dasar yang miring $\ang{3} = \qty{0.0524}{rad}$, $\alpha + \beta = 0.352/3.12 =
\qty{0.113}{rad} = \ang{6.5}$, maka $\alpha = \ang{3.5}$.
\end{solution}

\begin{solution}{exo:b3:orogenic-wedges:3}
Kedalaman $= (T - 10)/30$: apatit (U--Th)/He 1.3 sampai \qty{2.3}{km}; jejak fisi apatit
2.7 sampai \qty{3.7}{km}; jejak fisi zirkon $270/30 = \qty{9.0}{km}$. Pada
kedalaman \qty{5}{km} batuan berada pada suhu $10 + 150 = \qty{160}{\celsius}$: zirkon
(U--Th)/He, yang menutup pada 160 sampai \qty{200}{\celsius}, pada kedalaman 5.0 sampai
\qty{6.3}{km}.
\end{solution}

\begin{solution}{exo:b3:orogenic-wedges:4}
Panjang horizontal $2/\tan\ang{30} = \qty{3.46}{km}$; sepanjang kemiringan
$2/\sin\ang{30} = \qty{4.0}{km}$. Lapisan terangkat sebesar tinggi \omterm{def:b3:orogenic-wedges:thrust-geometry}{tanjakan},
\qty{2}{km}. Titik potong harus menempuh seluruh \omterm{def:b3:orogenic-wedges:thrust-geometry}{tanjakan}: pergeseran sebesar
\qty{4.0}{km}.
\end{solution}

\begin{solution}{exo:b3:orogenic-wedges:5}
$\phi = \arctan1.03 = \ang{45.8}$, $K = 5.08$; $\beta = \qty{0.105}{rad}$.
Untuk $\lambda = 0.4$: $\alpha + \beta = (0.105 + 0.85\times0.6)/(1 +
5.08\times0.6) = 0.615/4.05 = \qty{0.152}{rad} = \ang{8.7}$, $\alpha =
\ang{2.7}$. Untuk $\lambda = 0.9$: $0.190/1.51 = \qty{0.126}{rad} = \ang{7.2}$,
$\alpha = \ang{1.2}$. Peningkatan tekanan fluida menurunkan \omterm{def:b3:orogenic-wedges:critical-taper}{taper kritis}:
baji menjadi superkritis, meluncur ke depan pada dasarnya tanpa
menebal, dan frontnya maju sampai akresi dan erosi mengembalikannya
ke taper yang baru.
\end{solution}

\begin{solution}{exo:b3:orogenic-wedges:6}
$S = A/r = 6.0/1.5 = \qty{4.0}{km}$. Dibaca sebagai \omterm{def:b3:orogenic-wedges:thrust-folds}{lipatan detasemen}, $S = A/h =
6.0/3.0 = \qty{2.0}{km}$: setengah nilai sebenarnya. Kedua pembacaan harus
dibedakan berdasarkan bentuk lipatan dan keberadaan \omterm{def:b3:orogenic-wedges:thrust-geometry}{tanjakan}.
\end{solution}

\begin{solution}{exo:b3:orogenic-wedges:7}
$K = 2.12$ dan $1 + K(1 - 0.4) = 2.27$. Persamaan memberikan $\mu_b(1 -
\lambda_b) = (\alpha + \beta)\times2.27 - \beta$, dengan $\beta =
\qty{0.0349}{rad}$. Jura: $0.0611\times2.27 - 0.0349 = 0.10$. Selatan:
$0.218\times2.27 - 0.0349 = 0.46$. Byerlee pada tekanan hidrostatik:
$0.85\times0.6 = 0.51$. Dasar di selatan mendekati hukum Byerlee; dasar
Jura hampir lima kali lebih lemah, seperti yang diharapkan pada décollement
dalam halit, yang gesekannya rendah, di bawah fluida yang tidak bertekanan berlebih.
\end{solution}

\begin{solution}{exo:b3:orogenic-wedges:8}
$\kappa = \qty{e-6}{m^2/s} = \qty{31.6}{km^2/Ma}$, dan $v =
\qty{1}{km/Ma}$: $\mathrm{Pe} = 30/31.6 = 0.951$. Faktor $0.951/(1 -
e^{-0.951}) = 1.55$ menaikkan gradien dari 30 menjadi \qty{46.5}{\celsius/km}.
Untuk \qty{65}{\celsius}: $1 - e^{-\mathrm{Pe}\,z/L} = (55/900)(1 - e^{-0.951})
= 0.0375$, maka $\mathrm{Pe}\,z/L = 0.0382$ dan $z = \qty{1.21}{km}$, dibandingkan dengan
$55/30 = \qty{1.83}{km}$ tanpa erosi.
\end{solution}

\begin{solution}{exo:b3:orogenic-wedges:9}
$\mathrm dh/\mathrm dt = \qty{0.8}{km}/\qty{1.2}{Ma} = \qty{0.67}{mm/yr}$.
Dengan $f = 0.5$, $v = 0.67\times(1 - 0.5) = \qty{0.33}{mm/yr}$.
\end{solution}

\begin{solution}{exo:b3:orogenic-wedges:10}
Dengan $\alpha = 0$ dan $u = 1 - \lambda$: $\beta(1 + Ku) = \beta + \mu_b u$,
maka $\beta Ku = \mu_b u$ dan $\beta^* = \mu_b/K$, berapa pun $u$. Untuk Taiwan,
$0.85/5.08 = \qty{0.167}{rad} = \ang{9.6}$. Pada dasar yang lebih curam, kemiringan
permukaan kritis bernilai negatif: setiap baji dengan permukaan datar atau miring ke arah depan-pegunungan
sudah superkritis, dan ia meluncur pada dasarnya tanpa menebal. Sabuk
sesar dorong tidak dapat mencapai bentuk normalnya di sana.
\end{solution}

\begin{solution}{exo:b3:orogenic-wedges:11}
Ketika $\Lambda \to \infty$ dan $v \to 0$, $m \to 0$, maka $e^{-mz_c} \to 1$, dan
geotermnya lurus, $G_0 = G(z_c) = G$: $f \to (G - \Gamma)/G = 1 -
\Gamma/G = 1 - 6.5/30 = 0.78$. Kemiringannya kemudian $v/(1 - f) = vG/\Gamma =
4.6v$. Untuk $v = \qty{1}{mm/yr}$, $v/2\kappa = \qty{0.0158}{km^{-1}}$. Pada kedalaman
\qty{5}{km}: $m = 0.0158 + \sqrt{0.0158^2 + 1.257^2} = \qty{1.273}{km^{-1}}$,
$e^{-1.273\times1.21} = 0.215$, $f = 40.0\times0.215/44.7 = 0.19$. Pada kedalaman
\qty{20}{km}: $m = 0.0158 + 0.3146 = \qty{0.330}{km^{-1}}$, $e^{-0.399} =
0.671$, $f = 40.0\times0.671/44.7 = 0.60$.
\end{solution}

\begin{solution}{exo:b3:orogenic-wedges:12}
Penampang berupa segitiga dengan alas $W$ dan tinggi $W\tan(\alpha + \beta)$:
$A = W^2\tan(\alpha + \beta)/2$. Dengan $\mathrm dA/\mathrm dt = F_a$ dan $A =
0$ pada $t = 0$, $A = F_at$, yang memberikan $W$. Dengan $\tan\ang{6} = 0.105$:
$W = \sqrt{2\times40\times2/0.105} = \qty{39}{km}$ setelah \qty{2}{Ma}; $t =
60^2\times0.105/(2\times40) = \qty{4.7}{Ma}$ untuk mencapai \qty{60}{km}.
\end{solution}

\begin{solution}{pb:b3:orogenic-wedges:1}
\textbf{1.} $4.5 - 1.5 = \qty{3.0}{km}$ tinggi; $3.0/\tan\ang{25} =
\qty{6.43}{km}$ panjang horizontal; $3.0/\sin\ang{25} = \qty{7.10}{km}$ sepanjang
kemiringan.

\textbf{2.} Berdasarkan tinggi \omterm{def:b3:orogenic-wedges:thrust-geometry}{tanjakan}, \qty{3.0}{km}.

\textbf{3.} $S = A/r = 16.5/3.0 = \qty{5.5}{km}$.

\textbf{4.} Nilai yang sama, \qty{16.5}{km^2}: pada \omterm{def:b3:orogenic-wedges:thrust-folds}{lipatan tekukan-sesar}, luas kelebihan
setiap lapisan di atas \omterm{def:b3:orogenic-wedges:thrust-geometry}{bagian datar} atas adalah $Sr$. Sebagai \omterm{def:b3:orogenic-wedges:thrust-folds}{lipatan detasemen}, $16.5/4.0 =
\qty{4.1}{km}$: kurang seperempat.

\textbf{5.} $4\times(6.0 - 2.5) = \qty{14.0}{km}$, dua setengah kali
\qty{5.5}{km} pergeseran pada \omterm{def:b3:orogenic-wedges:thrust-geometry}{tanjakan} frontal.

\textbf{6.} Kuda frontal, yang paling dekat dengan depan-pegunungan: setiap \omterm{def:b3:orogenic-wedges:thrust-geometry}{tanjakan} baru
terbentuk di dinding bawah di depan yang terakhir. Setiap kuda, terangkat naik sepanjang \omterm{def:b3:orogenic-wedges:thrust-geometry}{tanjakannya},
mengangkat sesar dorong atap dan batuan penutup di atasnya seperti \omterm{def:b3:orogenic-wedges:thrust-folds}{lipatan tekukan-sesar} mengangkat
dinding gantungnya, sehingga atap terlipat di atas tumpukan.

\textbf{7.} $\alpha = \arctan(2.1/60) = \ang{2.00}$; $\beta =
\arctan(4.2/60) = \ang{4.00}$; taper \ang{6.0}.

\textbf{8.} $2.4 + 8.2 = \qty{10.6}{km}$, kurang dari sekitar \qty{15}{km}: dasarnya
masih bersifat gesekan, dan teori Coulomb berlaku.

\textbf{9.} $\phi = \ang{45.8}$, $\sin\phi = 0.717$, $K = 5.08$. Kering:
$\alpha + \beta = (0.0699 + 0.85)/6.08 = \qty{0.151}{rad} = \ang{8.7}$,
$\alpha = \ang{4.7}$.

\textbf{10.} $\lambda = 1000/2500 = 0.40$: $(0.0699 + 0.51)/4.05 =
\qty{0.143}{rad} = \ang{8.2}$, $\alpha = \ang{4.2}$. Sudut yang teramati \ang{2.0}
akan membuat baji sangat subkritis, menebal alih-alih maju:
fluida harus bertekanan berlebih.

\textbf{11.} Dengan $T = \qty{0.1049}{rad}$ dan $u = 1 - \lambda$, $T(1 + Ku) =
\beta + \mu_bu$ menghasilkan $u = (\beta - T)/(TK - \mu_b) = (0.0699 -
0.1049)/(0.5327 - 0.85) = 0.110$: $\lambda = 0.89$.

\textbf{12.} $\sigma_v = 2500\times9.81\times10\,600 = \qty{260}{MPa}$; $P =
0.89\times260 = \qty{231}{MPa}$; hidrostatik $1000\times9.81\times10\,600 =
\qty{104}{MPa}$. Fluida membawa \qty{127}{MPa} lebih besar daripada
kolom air.

\textbf{13.} $(0.0699 + 0.85\times0.3)/(1 + 5.08\times0.3) = 0.325/2.52 =
\qty{0.129}{rad} = \ang{7.4}$, $\alpha = \ang{3.4}$. Baji, pada \ang{2.0},
kini subkritis: ia berhenti maju dan menebal oleh \omterm{def:b3:orogenic-wedges:thrust-geometry}{sesar dorong luar-urutan}
di belakang frontnya hingga kemiringan permukaannya mencapai \ang{3.4}.

\textbf{14.} $F_a = \qty{10}{km/Ma}\times\qty{4.0}{km} = \qty{40}{km^2/Ma}$,
yaitu \qty{4.0e4}{m^3/yr} per kilometer sabuk.

\textbf{15.} $\bar e = 40/60 = \qty{0.67}{km/Ma} = \qty{0.67}{mm/yr}$.

\textbf{16.} $40\times100 = \qty{4000}{km^3/Ma} = \qty{4.0e6}{m^3/yr}$, atau
$\qty{1.0e10}{kg/yr}$ (10 juta ton per tahun).

\textbf{17.} $\mathrm dW/\mathrm dt = F_a/(W\tan(\alpha + \beta)) =
40/(60\times0.105) = \qty{6.3}{km/Ma}$.

\textbf{18.} Tidak: pada keadaan tunak erosi harus tetap memindahkan seluruh
fluks masuk, sehingga laju rata-rata tetap \qty{0.67}{mm/yr}. Namun batuan yang ditambahkan
di bawah bagian belakang terangkat di sana: erosi dan ekshumasi akan paling cepat
di atas bagian belakang, tempat umur pendinginan termuda akan ditemukan.

\textbf{19.} $v = \qty{0.667}{km/Ma}$: $\mathrm{Pe} = 0.667\times30/31.6 =
0.634$; $G_0 = 30\times0.634/(1 - e^{-0.634}) = 30\times1.35 =
\qty{40.5}{\celsius/km}$.

\textbf{20.} $1 - e^{-\mathrm{Pe}\,z/L} = (55/900)(1 - e^{-0.634}) = 0.0287$,
sehingga $\mathrm{Pe}\,z/L = 0.0291$ dan $z_c = \qty{1.38}{km}$; $G(z_c) =
40.5\,e^{-0.0291} = \qty{39.3}{\celsius/km}$.

\textbf{21.} $1.38/0.667 = \qty{2.07}{Ma}$, umur sampel pada
\qty{2.00}{km}: pada elevasi rata-rata, defleksi isoterm adalah nol.

\textbf{22.} $v/2\kappa = 0.667/63.1 = \qty{0.0106}{km^{-1}}$, $2\pi/\Lambda =
\qty{0.314}{km^{-1}}$: $m = 0.0106 + 0.3143 = \qty{0.325}{km^{-1}}$;
$e^{-0.325\times1.38} = 0.639$; $f = (40.5 - 6.5)\times0.639/39.3 = 0.552$.

\textbf{23.} Elevasi rata-rata adalah \qty{2.00}{km} dan umur rata-rata
\qty{2.068}{Ma}; $\sum\Delta t\,\Delta h = 0.823$ dan $\sum\Delta t^2 =
0.552$, sehingga kemiringannya \qty{1.49}{km/Ma}: laju semu sebesar
\qty{1.49}{mm/yr}.

\textbf{24.} $v = 1.49\times(1 - 0.552) = \qty{0.667}{mm/yr}$.

\textbf{25.} \omterm{def:b3:orogenic-wedges:thermochronology}{Termokronometer} dan keseimbangan fluks sepakat: baji berada
pada keadaan tunak, dan mengambil kemiringan semu apa adanya akan
melebih-lebihkan laju lebih dari dua kali lipat. Laju ekshumasi baji
Kestrel adalah \textbf{\qty[text-series-to-math=true]{0.67}{mm/yr}}.
\end{solution}
